% <<<PDFPAGE 1>>>
\markboth{}{}
\thispagestyle{book}
\vspace*{45pt}

{\color{bookblue}\hrule height 2.5pt}
\begin{center}
  \vspace{12pt}
  {\Huge\heiti\bfseries\color{bookblue} 大语言模型高效推理\par}
  \vspace{5pt}
  {\Large\heiti\bfseries\color{boxgray} 系统、算法与生产工程\par}
  \vspace{3pt}
  {\LARGE\heiti\bfseries\color{bookteal} 面试速查笔记\par}
  \vspace{10pt}
\end{center}
{\color{bookblue}\hrule height .9pt}

\vspace{35pt}
\begin{center}
\begin{tikzpicture}[
  x=25mm,y=9mm,
  neuron/.style={circle,draw,thick,minimum size=19pt,inner sep=0pt},
  connection/.style={gray!43,line width=.35pt},
  layer/.style={font=\heiti\bfseries\small}
]
  \foreach \i in {1,...,5} {
    \coordinate (input-\i) at (0,3-\i);
  }
  \foreach \i in {1,...,6} {
    \coordinate (hidden-\i) at (1,3.5-\i);
  }
  \foreach \i in {1,...,3} {
    \coordinate (output-\i) at (2,2-\i);
  }

  \foreach \i in {1,...,5} {
    \foreach \j in {1,...,6} {
      \draw[connection] (input-\i) -- (hidden-\j);
    }
  }
  \foreach \i in {1,...,6} {
    \foreach \j in {1,...,3} {
      \draw[connection] (hidden-\i) -- (output-\j);
    }
  }

  \foreach \i in {1,...,5} {
    \node[neuron,draw=bookblue,fill=bookblue!10] at (input-\i) {};
  }
  \foreach \i in {1,...,6} {
    \node[neuron,draw=bookteal,fill=bookteal!10] at (hidden-\i) {};
  }
  \foreach \i in {1,...,3} {
    \node[neuron,draw=orange,fill=orange!12] at (output-\i) {};
  }
  \node[layer,bookblue] at (0,2.65) {输入层};
  \node[layer,bookteal] at (1,3.5) {隐藏层};
  \node[layer,boxamber] at (2,1.65) {输出层};
\end{tikzpicture}

\vspace{25pt}
{\Large\heiti\bfseries\color{bookblue} 作者\par}
\smallskip
{\large\bfseries AI Engineering Insider\par}
Lamhot Siagian\par
\smallskip
{\small\color{bookteal}\url{https://www.linkedin.com/in/lamhotsiagian}\par}

\vspace{17pt}
2026 年第二版
\end{center}

\vfill
{\color{bookblue}\hrule height 1.5pt}
\smallskip
{\centering\small\itshape
涵盖量化、推测解码、KV 缓存、FlashAttention、vLLM、混合专家、\\
边缘推理、生产系统及推理时计算扩展\par}

\clearpage
% <<<PDFPAGE 2–4>>> 原书目录由主文件按中文页码生成。
